% !TeX program = xelatex
% Poisson bracket as the rate of change along Hamiltonian flow.
% Compile: xelatex poisson_evolution_tikz.tex
\documentclass[tikz,border=6pt]{standalone}
\usepackage{amsmath}
\usepackage[UTF8,fontset=fandol]{ctex}
\usepackage{fix-cm}
\usetikzlibrary{arrows.meta}

\begin{document}
\begin{tikzpicture}[>=Stealth,thick,every node/.style={font=\fontsize{15}{18}\selectfont}]
  % Non-conserved observable.
  \begin{scope}
    \node[blue!70!black] at (0,3.6) {$[f,H]\ne0$};
    \draw[blue!70!black] (0,0) ellipse (3.2 and 2.0);
    \begin{scope}
      \clip (0,0) ellipse (3.2 and 2.0);
      \foreach \c in {-2.0,-.7,.6,1.9} {
        \draw[red!65] (-3.4,{.45*(-3.4)+\c}) -- (3.4,{.45*3.4+\c});
      }
    \end{scope}
    \coordinate (A) at ({3.2*cos(135)},{2*sin(135)});
    \fill[black] (A) circle (2pt);
    \draw[green!45!black,->,line width=1.4pt] (A) -- ++(1.25,.78)
      node[above left] {$X_H$};
    \node[blue!70!black,anchor=west] at (2.5,-1.45) {$H=\text{常数}$};
    \node[red,anchor=west] at (2.15,2.5) {$f=\text{常数}$};
    \node[black!70,align=center] at (0,-2.75) {流线穿过 $f$ 的等值线\\因此 $f$ 随时间改变};
  \end{scope}

  % Conserved observable.
  \begin{scope}[xshift=8.3cm]
    \node[blue!70!black] at (0,3.6) {$[f,H]=0$};
    \draw[red!55,line width=3.5pt] (0,0) ellipse (3.2 and 2.0);
    \draw[blue!70!black,line width=1.2pt] (0,0) ellipse (3.2 and 2.0);
    \coordinate (B) at ({3.2*cos(135)},{2*sin(135)});
    \fill[black] (B) circle (2pt);
    \draw[green!45!black,->,line width=1.4pt] (B) -- ++(1.25,.78)
      node[above left] {$X_H$};
    \node[blue!70!black,anchor=west] at (2.2,-1.45) {$H=\text{常数}$};
    \node[red,anchor=west] at (2.05,1.7) {$f=\text{常数}$};
    \node[black!70,align=center] at (0,-2.75) {$f$ 的等值线与流线重合\\因此 $f$ 是守恒量};
  \end{scope}

  \node[green!40!black] at (4.15,4.45)
    {$\displaystyle \frac{df}{dt}=[f,H]\quad(f\text{ 不显含 }t)$};
\end{tikzpicture}
\end{document}
